\documentclass[11pt]{article}
\usepackage[margin=1in]{geometry}
\usepackage{amsmath,amssymb,amsthm}
\usepackage{booktabs}
\usepackage[hidelinks]{hyperref}
\usepackage{microtype}

\newtheorem{theorem}{Theorem}
\newtheorem{lemma}[theorem]{Lemma}
\newtheorem{corollary}[theorem]{Corollary}
\newtheorem{proposition}[theorem]{Proposition}
\theoremstyle{definition}
\newtheorem{definition}[theorem]{Definition}
\newtheorem{remark}[theorem]{Remark}

\newcommand{\Kgen}{K_{\mathrm{gen}}}
\newcommand{\Nk}{N_k}
\newcommand{\intr}{\operatorname{int}}

\title{Three Kobon numbers, a crossing-refined capacity theorem,\\
and the exact status of $K(14)$}
\author{Working draft --- author to be supplied}
\date{21 August 2026}

\begin{document}
\maketitle

\begin{abstract}
The Kobon triangle problem asks for the maximum number of non-overlapping
triangles determined by $n$ straight lines. We first separate three
quantities that the literature conflates: $a_3^s(n)$, triangular faces of
\emph{simple} arrangements; $K(n)$, triangular faces of \emph{arbitrary}
arrangements (the object of Clément--Bader's definition, in which a counted
triangle is ``realized by $3$ straight line segments''); and $\Kgen(n)$, the
maximum size of a family of triangles with pairwise disjoint interiors in
which other lines \emph{may} cross a counted triangle. Then
$a_3^s\le K\le\Kgen$, and we show that the classical segment-counting
arguments bound only the first two.

We prove a capacity inequality valid for $\Kgen$ with arbitrary parallels and
multiple points,
\[
 3|S|+C+2Q+\sum_p k_p(k_p-4)\ \le\ n(n-2),
\]
where $C$ counts line/triangle interior incidences, together with an exact
bounded-face budget; combining them gives the closed form
$\Kgen(n)\le\lfloor (n-2)(5n-3)/12\rfloor$, in which triple points are exactly
neutral. To our knowledge these are the only bounds proved for $\Kgen$.

We then use exact SAT certificates to settle small cases in this strongest
convention: $\Kgen(n)$ for every $n\le10$ --- with DRAT-verified
refutations at $n=4,\dots,9$ --- which are strictly stronger statements than
the corresponding classical ones because they also exclude crossed
families. On the constructive side we independently
verify, in exact rational arithmetic and with three implementations, a
$14$-line arrangement with $54$ pairwise interior-disjoint triangles, so
$\Kgen(14)\ge54$; and we prove that every currently known extremal
arrangement at $n=11,14,18,20$ attains its published count as its exact
optimum, so crossings never help inside known configurations. Finally we
observe that the entry $a(14)=54$ recorded as exact in OEIS A006066 on
21 August 2026 pairs a non-simple construction with a simple-arrangement
upper bound; the rigorous status is $54\le K(14)\le55$, and we describe the
decision procedure that closes it, together with two equality-branch
obstruction theorems ($n=14$ and $n=18$, machine-checked) and a proof that
the method cannot close the analogous branch at $n=20$.
\end{abstract}

\section{Three quantities}

Let $\mathcal A$ be an arrangement of $n\ge3$ distinct lines in the real
affine plane, not all parallel (\emph{essential}). A \emph{segment} of
$\mathcal A$ is a bounded maximal piece of a line between two consecutive
crossing points.

\begin{definition}\label{def:three}
\begin{enumerate}
\item[(i)] $a_3^s(n)$: the maximum number of triangular bounded faces of a
\emph{simple} arrangement of $n$ lines (no two parallel, no three
concurrent).
\item[(ii)] $K(n)$: the maximum, over all arrangements of $n$ lines, of the
number of triangles each of whose three sides is a single segment of
$\mathcal A$, no two of them overlapping. Equivalently --- since a line
entering the interior of a triangle must cross two of its sides and thereby
subdivide them --- the maximum number of triangular bounded faces of an
arbitrary arrangement.
\item[(iii)] $\Kgen(n)$: the maximum, over all arrangements, of the size of a
family $S$ of triangles bounded by triples of lines of $\mathcal A$ whose
open interiors are pairwise disjoint. Members of $S$ may be crossed by other
lines of $\mathcal A$.
\end{enumerate}
Clearly $a_3^s(n)\le K(n)\le \Kgen(n)$.
\end{definition}

Clément and Bader \cite{CB2007} define a Kobon triangle as ``a triangle that
is realized by 3 straight lines segments and which does not overlap with
other triangles'', i.e.\ they work with $K$; their Lemma 1 argument
enumerates the cases ``side shared by two triangles'' (which forces
concurrency at an endpoint) and ``segment used by no triangle'', and their
mod-$6$ improvement passes through \emph{perfect configurations}, defined as
arrangements of \emph{pairwise intersecting} lines in which every segment is
the side of exactly one triangle, and shown in their Lemma 2 to have no three
concurrent lines. Their bound
\begin{equation}\label{eq:CB}
 K(n)\ \le\ \lfloor n(n-2)/3\rfloor - \mathbf 1_{n \bmod 6\in\{0,2\}}
\end{equation}
is therefore a statement about $K$; the case of a counted triangle whose
sides are subdivided by crossings is not part of their case analysis, and
their claim that a segment serving two triangles forces a multiple point at
one of its endpoints fails once subdivided sides are admitted. Bartholdi,
Blanc and Loisel \cite{BBL} and Blanc \cite{Blanc} bound $a_3^s$. Thus:

\begin{quote}
\emph{The base segment bound $\lfloor n(n-2)/3\rfloor$ is established for
$K$ including non-simple arrangements by Felsner and Kriegel
\cite{FK1999} (so used, e.g., in \cite{PU1}); the mod-$6$ improvement
\eqref{eq:CB} for $K$ rests on the unpublished draft \cite{CB2007}; the
sharpest numerical bounds \cite{BBL,Blanc} are proved for $a_3^s$ only and
hence do not apply to non-simple record arrangements; and no published
bound is proved for $\Kgen$.}
\end{quote}

This has an immediate consequence for the current record. On 21 August 2026
OEIS A006066 \cite{OEIS} was updated to list $a(14)=54$ as exact, attributed
to Maiorana's arrangement \cite{Maiorana}. That arrangement is \emph{not}
simple: it has two triple points (on lines $\{0,4,10\}$ and $\{0,11,12\}$,
sharing line~$0$) and no parallels. The matching upper bound $54$ in the same
table is the BBL value $\lfloor n(n-7/3)/3\rfloor$, proved for simple
arrangements. For $K$, the peer-reviewed non-simple segment bound
\cite{FK1999} gives $56$ at $n=14$, and the mod-$6$ improvement
\eqref{eq:CB} of the draft \cite{CB2007} gives $55$. Hence:

\begin{proposition}\label{prop:status14}
With the published literature, $54\le K(14)\le 56$ unconditionally, and
$K(14)\le55$ if the draft \cite{CB2007} is granted; for the strongest
convention $54\le\Kgen(14)\le67$. In every case $a(14)=54$ is a lower bound
paired with a simple-arrangement upper bound, so its exactness is not
established.
\end{proposition}

\paragraph{Contributions.}
\begin{enumerate}
\item A capacity theorem for $\Kgen$ (Theorem~\ref{thm:capacity}), with exact
degeneracy penalty and crossing charge, and a closed-form ceiling
(Theorem~\ref{thm:ceiling}).
\item A certified initial segment in the strongest convention:
$\Kgen(4),\dots,\Kgen(9)=2,5,7,11,15,21$, each with a
\texttt{drat-trim}-verified refutation of $\Kgen(n)+1$, plus
$\Kgen(10)=25$ from a prior cube cover (Theorem~\ref{thm:ladder}).
\item An independently verified $\Kgen(14)\ge54$ (Theorem~\ref{thm:54}) and
an exact optimality census of all known extremal arrangements
(Theorem~\ref{thm:census}), which shows $K=\Kgen$ on each of them.
\item Two equality-branch obstruction theorems with machine-checked finite
certificates (Theorems~\ref{thm:M},~\ref{thm:M18}) and a proof that the same
method cannot close the analogous branch at $n=20$
(Proposition~\ref{prop:n20}).
\item The status audit of $K(14)$ (Proposition~\ref{prop:status14}) and the
decision procedure that closes it.
\end{enumerate}

\section{Notation}

For $\ell\in\mathcal A$ let $q_\ell$ be the number of lines parallel to
$\ell$ and $Q=\frac12\sum_\ell q_\ell$. A \emph{multipoint} has multiplicity
$k_p\ge3$; $\Nk$ counts points of multiplicity $k$. On $\ell$, $r_\ell$ is
the number of distinct crossings, $a_\ell$ the number of multipoints, and
$m_\ell=r_\ell-1$ the number of bounded elementary gaps. For an admissible
family $S$ (Definition~\ref{def:three}(iii)) set
\[
 \sigma_\ell=\#\{T\in S:\ \ell\text{ supports a side of }T\},\qquad
 \gamma_\ell=\#\{T\in S:\ \ell\cap\intr(T)\ne\emptyset\},
\]
so $\sum_\ell\sigma_\ell=3|S|$ and $C:=\sum_\ell\gamma_\ell$ is the total
number of line/triangle interior incidences. Write $T=|S|$ and
$\Delta=n(n-2)-3T$.

\section{The capacity theorem}

\begin{lemma}[incidence identity]\label{lem:identity}
$m_\ell+2a_\ell = n-2-q_\ell-\sum_{p\ni\ell}(k_p-4)$ for every
$\ell\in\mathcal A$.
\end{lemma}

\begin{proof}
The $n-1-q_\ell$ lines meeting $\ell$ do so in $r_\ell$ distinct points, and a
multipoint $p$ on $\ell$ absorbs $k_p-1$ of them into one point, so
$n-1-q_\ell=r_\ell+\sum_{p\ni\ell}(k_p-2)$. Hence
$m_\ell+2a_\ell=r_\ell-1+2a_\ell=n-2-q_\ell-\sum_{p\ni\ell}(k_p-2)+2a_\ell$,
and $a_\ell$ is the number of summands.
\end{proof}

\begin{lemma}[per-line capacity]\label{lem:perline}
$\sigma_\ell+\gamma_\ell\le m_\ell+2a_\ell$ for every $\ell\in\mathcal A$ and
every admissible $S$.
\end{lemma}

\begin{proof}[Proof sketch]
Each $T\in S$ with a side on $\ell$ determines the side interval
$I_T\subset\ell$; each $T$ crossed by $\ell$ determines the chord interval
$J_T=\ell\cap\intr(T)$. Disjointness of interiors makes the $J_T$ pairwise
disjoint, makes each $J_T$ disjoint from the interiors of side intervals on
the same side of $\ell$, and makes side intervals on a fixed side of $\ell$
interior-disjoint. Every such interval contains at least one elementary gap,
and two intervals can claim the same gap only when they lie on opposite sides
of $\ell$ and abut a common multipoint of $\ell$, each of which admits at
most two coincidences. Charging one gap per interval and two tokens per
multipoint yields the claim; the exhaustive case analysis is carried out in
the companion document \cite[\S3]{companion} and was re-derived by an
adversarial audit over $3{,}016$ arrangements and $812{,}486$ admissible
families with $5{,}900{,}863$ exact linewise checks.
\end{proof}

\begin{theorem}[capacity]\label{thm:capacity}
For every essential arrangement of $n$ distinct affine lines and every
admissible family $S$,
\begin{equation}\label{eq:capacity}
 3|S|+C+2Q+\sum_p k_p(k_p-4)\ \le\ n(n-2).
\end{equation}
\end{theorem}

\begin{proof}
Sum Lemma~\ref{lem:perline} over all lines and substitute
Lemma~\ref{lem:identity}; the left side gives $3|S|+C$, and on the right
$\sum_\ell q_\ell=2Q$ while each multipoint lies on exactly $k_p$ lines, so
$\sum_\ell\sum_{p\ni\ell}(k_p-4)=\sum_p k_p(k_p-4)$.
\end{proof}

Setting $C=0$ recovers, with explicit degeneracy corrections, the classical
segment bound for $K$; the point of \eqref{eq:capacity} is that the crossing
charge makes it valid for $\Kgen$.

\begin{remark}
Neither correction term can be replaced by the naive alternatives: the
variants with $\binom{k_p-2}{2}$ or $(k_p-2)^2$ in place of $k_p(k_p-4)$ are
false, with exact counterexamples in \cite[\S6]{companion}. The all-parallel
class is the unique inessential exception.
\end{remark}

\begin{corollary}\label{cor:tamura}
If $\mathcal A$ has no triple point then
$|S|\le\frac13\big(n(n-2)-C-2Q-\sum_{k\ge4}\Nk k(k-4)\big)\le n(n-2)/3$. In
general
$|S|\le n(n-2)/3+N_3-\frac13\big(C+2Q+\sum_{k\ge5}\Nk k(k-4)\big)$, so triple
points are the only degeneracy that can raise the ceiling, by at most one
triangle each.
\end{corollary}

\section{A closed-form ceiling for $\Kgen$}

\begin{lemma}[face budget]\label{lem:face}
$|S|+C+Q+\sum_p\binom{k_p-1}{2}\le\binom{n-1}{2}$.
\end{lemma}

\begin{proof}
$\mathcal A$ has exactly $\binom{n-1}{2}-Q-\sum_p\binom{k_p-1}{2}$ bounded
faces. The interior of $T\in S$ contains at least $1+c_T$ bounded faces,
where $c_T$ is the number of lines crossing it, and distinct members have
disjoint interiors, so they consume disjoint face sets; summing gives
$|S|+C$.
\end{proof}

\begin{theorem}[closed-form ceiling]\label{thm:ceiling}
$\Kgen(n)\le\big\lfloor (n-2)(5n-3)/12\big\rfloor$ for all $n\ge3$. In the
combined inequality used, multiplicity $3$ has coefficient exactly $0$ and
every $k\ge4$ has coefficient $k(k-4)+3\binom{k-1}{2}>0$.
\end{theorem}

\begin{proof}
Add \eqref{eq:capacity} to three times Lemma~\ref{lem:face}:
\[
 6|S|+4C+5Q+\sum_{k\ge3}\Nk\Big(k(k-4)+3\binom{k-1}{2}\Big)\le
 n(n-2)+3\binom{n-1}{2}=\frac{(n-2)(5n-3)}{2}.
\]
The bracket is $0,9,23,42,\dots$ for $k=3,4,5,6$, hence nonnegative.
\end{proof}

\begin{lemma}[reduction at zero slack]\label{lem:reduction}
If $\Delta=0$ and $\mathcal A$ has no triple point, then $C=0$, $Q=0$, every
multipoint has $k_p=4$, and every member of $S$ is a bounded triangular face;
in particular $|S|\le K(n)$ is attained by faces.
\end{lemma}

\begin{proof}
\eqref{eq:capacity} forces $C+2Q+\sum_{k\ge4}\Nk k(k-4)\le0$ with all terms
nonnegative.
\end{proof}

Among known values, \eqref{eq:capacity} is attained exactly at
$n=3,5,9,15,17$ --- precisely the $n$ where the Tamura value $n(n-2)/3$ is
realised --- and there all three quantities of Definition~\ref{def:three}
coincide.

\section{Certified small values in the strongest convention}

The decision problem ``is there an $n$-line arrangement with an admissible
family of size $T$?'' is encoded in propositional logic over slope-sorted
combinatorial data: variables for crossing versus parallel, concurrency of
triples, the order of crossings along each line, the selection, and
triangle/line crossing incidences, with the per-line inequalities of
Lemma~\ref{lem:perline} and the budget of Lemma~\ref{lem:face} as exact
cardinality constraints. All degeneracies are free, so a verdict covers
crossed and degenerate configurations simultaneously.

\begin{theorem}\label{thm:ladder}
In the convention of Definition~\ref{def:three}(iii),
\[
 \Kgen(4),\dots,\Kgen(9)\;=\;2,\,5,\,7,\,11,\,15,\,21,
\]
and each value is certified: the instance $(n,\Kgen(n)+1)$ is unsatisfiable
with a \texttt{drat-trim}-verified proof (Table~\ref{tab:ladder}), and the
matching lower bounds are the classical constructions. Together with the
prior eleven-cube cover for $n=10$ this gives $\Kgen(n)$ for all $n\le10$.
\end{theorem}

\begin{table}[t]
\centering
\begin{tabular}{rrrrrrr}
\toprule
$n$ & $\Kgen(n)$ & refuted $T$ & variables / clauses & core clauses &
core lemmas & verify (s)\\
\midrule
4 & 2 & 3 & $207$ / $994$ & $90$ & $131$ & $0.05$\\
5 & 5 & 6 & $972$ / $5{,}257$ & $424$ & $558$ & $0.04$\\
6 & 7 & 8 & $3{,}655$ / $20{,}681$ & $4{,}541$ & $7{,}018$ & $0.48$\\
7 & 11 & 12 & $10{,}682$ / $63{,}405$ & $12{,}255$ & $31{,}595$ & $4.47$\\
8 & 15 & 16 & $27{,}451$ / $165{,}367$ & $31{,}218$ & $411{,}278$ & $117.6$\\
9 & 21 & 22 & $61{,}569$ / $377{,}405$ & $66{,}346$ & $2{,}150{,}510$ & $1060.3$\\
\bottomrule
\end{tabular}
\caption{Certified initial segment. Every row is a \texttt{s VERIFIED}
\texttt{drat-trim} run on a kissat proof of the refutation of
$T=\Kgen(n)+1$; proofs range from $12$\,kB to $328.5$\,MB and the resolution
counts from $582$ to $8.7\cdot10^7$. Ledger:
\texttt{scratch/kobon/ladder\_certificates.json}.}
\label{tab:ladder}
\end{table}

These statements are strictly stronger than the corresponding classical ones,
since they also exclude crossed families: e.g.\ $\Kgen(8)=15$ rules out the
value $16=\lfloor 8\cdot 6/3\rfloor$ left open for $\Kgen$ by the segment
bound, and $\Kgen(6)=7$ rules out $8$.

\section{A verified lower certificate at $n=14$}

\begin{theorem}\label{thm:54}
$\Kgen(14)\ge54$, witnessed by an arrangement with exact rational
coefficients.
\end{theorem}

\begin{proof}
From the SHA-256-pinned source file of \cite{Maiorana} we verified: the $14$
lines are pairwise non-proportional; the computed multipoint census is exactly
two triple points sharing a line, with $Q=0$, matching the file's declared
data; exactly $54$ triples have interiors met by no other line, computed with
two independent predicates, the second of which also rejects the case of a
line through a vertex separating the opposite side; and the resulting family
passes the audited exact-rational checker \texttt{verify\_selection}, which
validates non-degeneracy and pairwise interior-disjointness in homogeneous
rational coordinates. The same verification succeeds for all fifteen
published solutions.
\end{proof}

\section{Exact optimality of the known witnesses}

For a \emph{fixed} arrangement, the largest admissible family is the
independence number of the overlap graph on non-degenerate triples, with
edges joining pairs of intersecting interiors. We compute it exactly:
interiors are compared in homogeneous rational arithmetic, and optimality is
certified by SAT (a family of size $K$ is a model, UNSAT at $K+1$ certifies
the optimum).

\begin{theorem}\label{thm:census}
For each of the following the exact admissible optimum equals the published
count, with certified optimality: the fifteen $14$-line $54$-triangle
arrangements of \cite{Maiorana}; two exact rational $53$-triangle $14$-line
reconstructions; an $11$-line $32$-triangle arrangement; an $18$-line
$93$-triangle reconstruction; and a $20$-line $116$-triangle reconstruction.
In particular $K=\Kgen$ on every known extremal arrangement.
\end{theorem}

Consequently any improvement on the records must come from new arrangements,
and by Corollary~\ref{cor:tamura} any improvement past the Tamura value must
use triple points --- which is exactly the mechanism of Maiorana's step from
$53$ to $54$.

\section{Equality-branch obstructions}

When $\Delta$ is small the admissible degeneracy patterns form a short list,
and saturation of Lemma~\ref{lem:perline} forces the multiplicities $m_i$ of
double-extreme vertices at consecutive direction classes to satisfy
\begin{equation}\label{eq:cycle}
 m_{i-1}+m_i=2s_i,\qquad 0\le m_i\le s_is_{i+1},\qquad m_i\in\mathbb Z,
\end{equation}
cyclically in the class sizes $(s_1,\dots,s_R)$. For odd $R$ the rational
solution is unique; for even $R$ solvability requires
$\sum_i(-1)^i 2s_i=0$.

\begin{theorem}[$n=14$]\label{thm:M}
No arrangement of $14$ lines with exactly three parallel pairs, or one
parallel triad, and with no three lines concurrent admits an admissible
family of $54$ triangles.
\end{theorem}

\begin{proof}
Paired layout: sizes are three $2$'s and eight $1$'s, $R=11$ odd, and the
unique solution of \eqref{eq:cycle} violates an interval constraint for each
of the $\binom{11}{3}=165$ placements. Triad layout: one $3$ and eleven
$1$'s, $R=12$ even, and the alternating sum is $\pm4\ne0$ at each of the $12$
placements. Both enumerations are certified by an exact integer program
(pure standard library, no floating point).
\end{proof}

\begin{theorem}[$n=18$]\label{thm:M18}
No arrangement of $18$ lines with exactly three parallel pairs, or one
parallel triad, and with no three lines concurrent admits an admissible
family of $94$ triangles.
\end{theorem}

\begin{proof}
Same mechanism with $R=15$ and $R=16$: all $\binom{15}{3}=455$ paired
placements violate an interval constraint (gap-parity classes $140+90+225$,
$19$ dihedral orbits with exact multiplicities, reflection quotients $231$
and $9$ all asserted by the certificate), and all $16$ triad placements have
alternating sum $\pm4$.
\end{proof}
\begin{proposition}[the method stops at $n=20$]\label{prop:n20}
At $n=20$, $T=117$ we have $\Delta=9$ and $2Q=6$, leaving three units of
slack, so \eqref{eq:cycle} must be relaxed by a defect vector, and feasible
defects exist: for the paired layout ($R=17$) a defect of $2$ at one doubled
class with run lengths $(u,0,14-u)$, e.g.\ $(6,0,8)$; for the triad layout
($R=18$) a defect of $4$ with all $m_i=1$. Hence no argument using only
\eqref{eq:cycle} can exclude $T=117$ at $n=20$.
\end{proposition}

\section{The status of $K(14)$, and how to close it}

Table~\ref{tab:windows} collects what is proved. The $n=14$ row is the
interesting one: the construction gives $54$, \eqref{eq:CB} gives $55$ for
$K$, and Theorem~\ref{thm:ceiling} gives $67$ for $\Kgen$. Two decision
instances close the gap in the strongest convention: $(14,55)$, and --- if
that is unsatisfiable --- nothing further is needed for $K$, while for
$\Kgen$ the rungs $T=56,\dots,67$ remain, of which $T=56$ is the first where
Lemma~\ref{lem:reduction} applies in the triple-point-free case. Both
$(14,55)$ and $(14,56)$ instances (about $1.34$--$1.35\cdot10^6$ variables
and $7.8\cdot10^6$ clauses) are running at the time of writing.

\begin{table}[t]
\centering
\begin{tabular}{rrrrl}
\toprule
$n$ & construction & $K(n)\le$ \eqref{eq:CB} & $\Kgen(n)\le$ Thm.~\ref{thm:ceiling} & certified here\\
\midrule
8 & 15 & 15 & 18 & $\Kgen(8)=15$ (DRAT)\\
9 & 21 & 21 & 24 & $\Kgen(9)=21$\\
10 & 25 & 26 & 31 & $\Kgen(10)=25$ (cube cover, prior)\\
11 & 32 & 33 & 39 & ---\\
12 & 38 & 39 & 47 & ---\\
13 & 47 & 47 & 56 & ---\\
14 & \textbf{54} (verified) & 55 & 67 & $\Kgen(14)\ge54$\\
18 & 93 & 95 & 116 & ---\\
20 & 116 & 119 & 145 & ---\\
\bottomrule
\end{tabular}
\caption{Constructions are admissible families, hence lower bounds for all
three quantities. Column three is the Clément--Bader draft bound for $K$
(faces in arbitrary arrangements); its base value $\lfloor n(n-2)/3\rfloor$
is peer-reviewed for non-simple arrangements \cite{FK1999}, while the $-1$
for $n\equiv0,2\pmod6$ rests on the draft \cite{CB2007}. Column four is
Theorem~\ref{thm:ceiling}, the only bound we know of that is proved for
$\Kgen$.}
\label{tab:windows}
\end{table}

\section{Open problems}

\begin{enumerate}
\item \textbf{Kill the triple-point credit.} The $2a_\ell$ tokens in
Lemma~\ref{lem:perline} are what allow $\sum_p k_p(k_p-4)<0$ and open the
ceiling above Tamura's value. Conjecture: they are never simultaneously
realisable, so $\Kgen(n)\le n(n-2)/3$ with degeneracies. Evidence: equality
in \eqref{eq:capacity} occurs only at $N_3=0$ in the entire known census, and
Theorem~\ref{thm:census} shows crossings never help inside known witnesses.
\item \textbf{Decide $(14,55)$}, closing $K(14)$ and hence the freshly
recorded OEIS entry, and then the rungs up to $67$ for $\Kgen$.
\item \textbf{Defect-tolerant obstructions.} Proposition~\ref{prop:n20}
localises the missing ingredient: a constraint on where a defect may sit.
\item \textbf{A broad mod-$6$ refinement}, i.e.\ derive the $-1$ correction
of \eqref{eq:CB} inside the framework of \eqref{eq:capacity}, making the
$K$-bound rigorous for arbitrary arrangements and removing the reliance on an
unpublished draft.
\item \textbf{Separation.} Is $K(n)<\Kgen(n)$ for some $n$? The first
candidates are $n=13$ ($T=48$ requires a triple point and is excluded for
$K$ by \eqref{eq:CB}) and $n=11$ ($T=33$); both decision instances are
running.
\end{enumerate}

\appendix
\section{Artefacts and replay}\label{app:artifacts}

Release bundle \texttt{math/kobon/release/kobon-2026-08/} with
\texttt{MANIFEST.sha256}; documentation hub \texttt{math/kobon/docs/}.

\begin{itemize}
\item \texttt{verification/maiorana\_verify.py} --- fail-closed certification
of Theorem~\ref{thm:54} (exit-$0$ transcript included).
\item \texttt{verification/sweep\_all\_verify.py} --- all fifteen witnesses.
\item \texttt{verification/max\_packing.py} --- exact overlap graph and SAT
optimality, used for Theorem~\ref{thm:census}.
\item \texttt{scratch/kobon/broad\_ladder.py} --- ladder prober used for
Theorem~\ref{thm:ladder}; instances \texttt{ladder\_n8\_t16.cnf},
\texttt{ladder\_n9\_t22.cnf} and the DRAT proofs.
\item \texttt{proofs/appendix\_independent\_check.py},
\texttt{proofs/appendix\_check\_n18.py} --- exact integer certificates for
Theorems~\ref{thm:M} and~\ref{thm:M18}.
\item \texttt{proofs/defect\_lemma.md} --- analysis behind
Proposition~\ref{prop:n20}.
\item \texttt{scripts/engine.py} --- exact rational geometry core and CNF
builder.
\end{itemize}

\begin{thebibliography}{9}
\bibitem{BBL} N.~Bartholdi, J.~Blanc, S.~Loisel,
\emph{On simple arrangements of lines and pseudo-lines in $\mathbb P^2$ and
$\mathbb R^2$ with the maximum number of triangles}, Contemp.\ Math.\
\textbf{453} (2008), 105--116; arXiv:0706.0723.
\bibitem{Blanc} J.~Blanc, \emph{The best polynomial bounds for the number of
triangles in a simple arrangement of $n$ pseudo-lines}, Geombinatorics
\textbf{21} (2011), no.~1, 5--14; preprint arXiv:0801.2845.
\bibitem{CB2007} G.~Clément, J.~Bader, \emph{Tighter upper bound for the
number of Kobon triangles}, draft, ETH Zürich, 21 December 2007.
\bibitem{FK1999} S.~Felsner, K.~Kriegel, \emph{Triangles in Euclidean
arrangements}, Discrete Comput.\ Geom.\ \textbf{22} (1999), no.~3,
429--438.
\bibitem{Savchuk} P.~Savchuk, \emph{Constructing optimal Kobon triangle
arrangements via table encoding, SAT solving, and heuristic straightening},
arXiv:2507.07951 (2025).
\bibitem{PU1} R.~Parpalak, D.~Utkin, \emph{Enumeration and classification of
triangle-maximal pseudoline arrangements}, arXiv:2607.29236 (2026).
\bibitem{PU2} R.~Parpalak, D.~Utkin, \emph{The $18\cdot2^t+1$
triangle-maximal series of straight lines}, arXiv:2604.22035 (2026).
\bibitem{Maiorana} A.~Maiorana, \emph{Kobon triangles, $k=14$}, dataset,
\url{https://github.com/rufio72/kobon_triangles_k14}, commit
\texttt{e47c7cfd}.
\bibitem{OEIS} OEIS Foundation Inc., \emph{Sequence A006066: Kobon
triangles}, \url{https://oeis.org/A006066} (revision \#227, 21 August 2026).
\bibitem{companion} Companion technical document
\texttt{math/kobon/paper\_capacity.md}: full case analysis of
Lemma~\ref{lem:perline}, audit protocols, refutations of weaker penalty
forms, and the complete experiment ledger.
\end{thebibliography}

\end{document}
